% !TEX encoding = UTF-8 Unicode
% !TEX root =  Main.tex
%------------------------------------------------------------------------------
%	Formatvorlage für wissenschaftliche Arbeiten (Diplomarbeit, Bachelorarbeit, Masterarbeit)
%------------------------------------------------------------------------------
%	Version: 2.0
%	Datum: 06.06.2025


% Dokumentenkopf ---------------------------------------------------------------------------------------
%   Diese Vorlage basiert auf "scrreprt" aus dem koma-script.
% ------------------------------------------------------------------------------------------------------
\documentclass[
    12pt, % Schriftgröße
    DIV=10, % Änderung der Größe des Satzspiegels (bedruckbarer Bereich einer Seite), nur in Verbindung mit koma-script verwendbar
    ngerman, % für Umlaute, Silbentrennung etc.
    a4paper, % Papierformat
    oneside, % einseitiges Dokument
    titlepage, % es wird eine Titelseite verwendet
    parskip=half, % Abstand zwischen Absätzen (halbe Zeile)
    headings=normal, % Größe der Überschriften verkleinern
    listof=totoc, % Verzeichnisse im Inhaltsverzeichnis aufführen
    bibliography=totoc, % Literaturverzeichnis im Inhaltsverzeichnis aufführen
    index=totoc, % Index im Inhaltsverzeichnis aufführen
    captions=tableheading, % Beschriftung von Tabellen unterhalb ausgeben
    final % Status des Dokuments (final/draft)
]{scrreprt}

% UTF8 und T1 Fontencoding -----------------------------------------------------------------------------
\usepackage[utf8]{inputenc}
\usepackage[T1]{fontenc}
\usepackage{tikz} % Ensure TikZ/PGF is loaded after fontenc/inputenc

% !TEX encoding = UTF-8 Unicode
% !TEX root =  Main.tex

% Meta-Informationen ------------------------------------------------------------------------------------
%   Definition von globalen Parametern, die im gesamten Dokument verwendet
%   werden können (z.B auf dem Deckblatt etc.).
%
%   ACHTUNG: Wenn die Texte Umlaute oder ein Esszet enthalten, muss der folgende
%            Befehl bereits an dieser Stelle aktiviert werden:
%            \usepackage[latin1]{inputenc}
% -------------------------------------------------------------------------------------------------------
\newcommand{\titel}{Konzeptionierung und Entwicklung einer Softwarelösung für einen automatisierten und
datenschutzkonformen Bearbeitungsprozess von Anträgen am Beispiel des AStA der HWR Berlin}
\newcommand{\untertitel}{}
\newcommand{\untertitelDeckblatt}{}
\newcommand{\art}{Hausarbeit}
\newcommand{\fachgebiet}{}
\newcommand{\autor}{Antonia Schikora}
\newcommand{\keywords}{Hausarbeit, Antonia Schikora}
\newcommand{\fachbereich}{}
\newcommand{\fachrichtung}{Wirtschaftsinformatik\xspace}
\newcommand{\studienjahrgang}{WI23A}
\newcommand{\semester}{4. Semester}
\newcommand{\matrikelnr}{tbd}
\newcommand{\erstgutachter}{Prof. Dr.-Ing. Sebastian Schlesinger}
\newcommand{\zweitgutachter}{Denny Heise}
\newcommand{\hochschule}{Hochschule für Wirtschaft und Recht Berlin}
\newcommand{\ort}{Berlin}
\newcommand{\ausbildungsbetrieb}{Deloitte GmbH} 
\newcommand{\logo}{Logo_HWR.png}
\newcommand{\creator}{LaTeX}
\newcommand{\datum}{\today}  % Fügt das aktuelle Datum ein


% !TEX encoding = UTF-8 Unicode
% !TEX root =  Main.tex

%%%%%%%%%%%%%%%%%%%%%%%%%%%%%%%%%%%%%%%%%%%%%%%%%%%%%%%%%%%%%%%%%%%%%%%
% GRUNDLEGENDE DOKUMENTEINSTELLUNGEN
%%%%%%%%%%%%%%%%%%%%%%%%%%%%%%%%%%%%%%%%%%%%%%%%%%%%%%%%%%%%%%%%%%%%%%%
% Seitenlayout und Kopfzeilen
\usepackage[
    automark,     % Kapitelangaben in Kopfzeile automatisch erstellen
    headsepline,  % Trennlinie unter Kopfzeile
    ilines        % Trennlinie linksbündig ausrichten
]{scrlayer-scrpage}

% Seitengeometrie und Abstände
\usepackage{geometry}
\usepackage{setspace}

% Fortlaufende Nummerierung
\usepackage{chngcntr}

%%%%%%%%%%%%%%%%%%%%%%%%%%%%%%%%%%%%%%%%%%%%%%%%%%%%%%%%%%%%%%%%%%%%%%%
% SPRACHE UND SCHRIFT
%%%%%%%%%%%%%%%%%%%%%%%%%%%%%%%%%%%%%%%%%%%%%%%%%%%%%%%%%%%%%%%%%%%%%%%
% Spracheinstellungen
\usepackage[ngerman]{babel}

% Schriftart und Kodierung
\usepackage{newtxtext,newtxmath}  % Times New Roman als Hauptschrift
\renewcommand{\familydefault}{\rmdefault}
\addtokomafont{disposition}{\rmfamily}
\addtokomafont{chapter}{\fontsize{13pt}{13pt}\selectfont}
\addtokomafont{section}{\fontsize{13pt}{13pt}\selectfont}
\addtokomafont{subsection}{\fontsize{13pt}{13pt}\selectfont}

% Schriftgrößen und Textauszeichnung
\usepackage{relsize}      % Relative Schriftgrößen
\usepackage[normalem]{ulem}  % Bessere Unterstreichungen
\usepackage{textcomp}     % Zusätzliche Symbole
\usepackage{microtype}

%%%%%%%%%%%%%%%%%%%%%%%%%%%%%%%%%%%%%%%%%%%%%%%%%%%%%%%%%%%%%%%%%%%%%%%
% ABBILDUNGEN UND GRAFIKEN
%%%%%%%%%%%%%%%%%%%%%%%%%%%%%%%%%%%%%%%%%%%%%%%%%%%%%%%%%%%%%%%%%%%%%%%
% Grundlegende Grafikunterstützung
\usepackage{graphicx}
\graphicspath{{Bilder/}}

% Erweiterte Grafikfunktionen
\usepackage{svg}
\usepackage{tikz}
\usepackage[vflt]{floatflt}  % Textumfluss
\usepackage{subfigure}       % Teilabbildungen
\usepackage{rotating}        % Rotation von Elementen

%%%%%%%%%%%%%%%%%%%%%%%%%%%%%%%%%%%%%%%%%%%%%%%%%%%%%%%%%%%%%%%%%%%%%%%
% HYPERLINKS UND PDF-EINSTELLUNGEN
%%%%%%%%%%%%%%%%%%%%%%%%%%%%%%%%%%%%%%%%%%%%%%%%%%%%%%%%%%%%%%%%%%%%%%%
% URL-Formatierung
\usepackage{url}

% PDF-Optionen und Hyperlinks
\usepackage[
    bookmarks,
    bookmarksopen=true,
    colorlinks=true,
    linkcolor=black,     % Interne Links (wie im Inhaltsverzeichnis) in Schwarz
    anchorcolor=black,
    citecolor=blue,
    filecolor=magenta,
    menucolor=black,
    urlcolor=cyan,
    plainpages=false,
    pdfpagelabels,
    hypertexnames=false,
    linktocpage,
    pdftitle={\titel\untertitel},
    pdfauthor={\autor},
    pdfcreator={\creator},
    pdfsubject={\titel\untertitel},
    pdfkeywords={\keywords}
]{hyperref}

%%%%%%%%%%%%%%%%%%%%%%%%%%%%%%%%%%%%%%%%%%%%%%%%%%%%%%%%%%%%%%%%%%%%%%%
% VERZEICHNISSE UND REFERENZEN
%%%%%%%%%%%%%%%%%%%%%%%%%%%%%%%%%%%%%%%%%%%%%%%%%%%%%%%%%%%%%%%%%%%%%%%
% Glossar und Abkürzungen
\usepackage[acronym,nomain,nopostdot,nonumberlist]{glossaries}
\makeglossaries%
\setglossarystyle{list}
\renewcommand*{\glsnamefont}[1]{\textrm{#1}}

% Literaturverzeichnis
% \usepackage[square,numbers,sort&compress]{natbib}
\usepackage{csquotes} % wurde im Log empfohlen
\usepackage[
    backend=biber,
    style=authoryear,
    language=ngerman,
    backref=true
]{biblatex}
\addbibresource{Literatur/Bibliographie.bib}


% To set square brackets for parencite, use:
\renewcommand*{\mkbibparens}[1]{\mkbibbrackets{#1}}
\DeclareRobustCommand{\addcomma}{,}
\DeclareLabelalphaTemplate{
  \labelelement{
    \field[strwidth=3,strside=left]{labelname}
  }
  \labelelement{
    \field[strwidth=2,strside=right]{year}
  }
}
\DeclareFieldFormat{labelyear}{%
  \iffieldundef{year}
    {o.\,J.} % ohne Jahr
    {#1}
}
\DeclareFieldFormat{labelname}{%
  \iffieldundef{name}
    {o.\,A.} % ohne Autor
    {#1}
}
\addbibresource{Literatur/Bibliographie.bib}
% Index
\usepackage{makeidx}
\usepackage{etoolbox}

%%%%%%%%%%%%%%%%%%%%%%%%%%%%%%%%%%%%%%%%%%%%%%%%%%%%%%%%%%%%%%%%%%%%%%%
% MATHEMATIK UND SYMBOLE
%%%%%%%%%%%%%%%%%%%%%%%%%%%%%%%%%%%%%%%%%%%%%%%%%%%%%%%%%%%%%%%%%%%%%%%
\usepackage{amsmath,amsfonts}  % Mathematische Symbole
\usepackage{eurosym}           % Währungssymbole

%%%%%%%%%%%%%%%%%%%%%%%%%%%%%%%%%%%%%%%%%%%%%%%%%%%%%%%%%%%%%%%%%%%%%%%
% TABELLEN UND AUFLISTUNGEN
%%%%%%%%%%%%%%%%%%%%%%%%%%%%%%%%%%%%%%%%%%%%%%%%%%%%%%%%%%%%%%%%%%%%%%%
% Tabellenformatierung
\usepackage{tabularx}
\usepackage{multirow}
\usepackage{array}
\usepackage{longtable}
\usepackage{ltxtable}

% Listen und Aufzählungen
\usepackage{paralist}
\usepackage{enumitem}
\setlist[description]{font=\rmfamily}

%%%%%%%%%%%%%%%%%%%%%%%%%%%%%%%%%%%%%%%%%%%%%%%%%%%%%%%%%%%%%%%%%%%%%%%
% PROGRAMMCODE UND LISTINGS
%%%%%%%%%%%%%%%%%%%%%%%%%%%%%%%%%%%%%%%%%%%%%%%%%%%%%%%%%%%%%%%%%%%%%%%
\usepackage{listings}
\usepackage{xcolor}

% Grundlegende Farben für Listings
\definecolor{hellgelb}{rgb}{1,1,0.9}
\definecolor{colKeys}{rgb}{0.8,0,0.5}
\definecolor{colIdentifier}{rgb}{0.6,0,0.3}
\definecolor{colComments}{rgb}{0,0.5,0}
\definecolor{colString}{rgb}{0,0,1}

% Listing-Standardeinstellungen
\lstset{
    float=htbp,
    basicstyle=\ttfamily\color{black}\small\smaller,
    identifierstyle=,%\color{colIdentifier},
    keywordstyle=\color{colKeys}\bfseries,
    stringstyle=\color{colString},
    commentstyle=\color{colComments},
    columns=flexible,
    tabsize=4,
    frame=single,
    extendedchars=true,
    showspaces=false,
    showstringspaces=false,
    numbers=left,
    numberstyle=\tiny,
    breaklines=true,
    backgroundcolor=\color{hellgelb},
    breakautoindent=true,
    escapeinside={(*}{*)},
    literate={Ö}{{\"O}}1 {Ä}{{\"A}}1 {Ü}{{\"U}}1 {ß}{{\ss}}2 {ü}{{\"u}}1 {ä}{{\"a}}1 {ö}{{\"o}}1 {µ}{\textmu}1
}

% Sprachdefinitionen für Listings
% CSS
\lstdefinelanguage{CSS}{
    morestring=[b]',
    morestring=[b]",
    comment=[l]{/*}{*/},
    sensitive=false,
    morekeywords={accelerator,adjust,after,align,attachment,azimuth,background,before,behavior,binding,border,bottom,bottomright,bottomleft,break,caption-side,char,clear,clip,color,colors,cue,cursor,collapse,decoration,direction,display,elevation,empty-cells,family,filter,float,flow,focus,font,grid,height,image,increment,indent,input,inside,ime-mode,include-source,justify,last,layer,left,letter,line,list,margin,marker,marks,max,min,mode,modify,moz,offsett,opacity,orphans,outline,overflow,overflow-X,overflow-Y,overhang,position,position-x,position-y,padding,page,pause,pitch,play-during,position,quotes,radius,range,repeat,replace,reset,richness,right,ruby,set-link-source,select,shadow,size,spacing,speak,speak-header,speak-numeral,speak-punctuation,speech-rate,stress,stretch,style,table-layout,text,transform,text-autospace,text-kashida-space,top,topleft,topright,type,underline,unicode-bidi,use-link-source,user,variant,vertical,visibility,voice-family,volume,white-space,weight,widows,width,word,wrap,word-wrap,writing-mode,z-index,zoom}
}

% JavaScript
\lstdefinelanguage{JavaScript}{
    keywords={typeof, new, true, false, catch, function, return, null, catch, switch, var, if, in, while, do, else, case, break},
    keywordstyle=\color{colKeys}\bfseries,
    ndkeywords={class, export, boolean, throw, implements, import, this},
    ndkeywordstyle=\color{darkgray}\bfseries,
    sensitive=false,
    comment=[l]{//},
    morecomment=[s]{/*}{*/},
    morestring=[b]',
    morestring=[b]"
}

% TypoScript
\lstdefinelanguage{TypoScript}{
    keywords=[1]{PAGE, HTML, TEXT, COA, COA\_INT, FILE, IMAGE, IMG\_RESOURCE, CLEARGIF, CONTENT, RECORDS, HMENU, CTABLE, OTABLE, COLUMNS, HRULER, IMGTEXT, CASE, LOAD\_REGISTER, RESTORE\_REGISTER, FORM, SEARCHRESULT, USER, USER\_INT, TEMPLATE, FLUIDTEMPLATE, MULTIMEDIA, SVG, EDITPANEL, GIFBUILDER, GMENU, TMENU, TMENUITEM, IMGMENU, IMGMENUITEM, JSMENU, JSMENUITEM, BOX},
    keywordstyle=[1]{\color{blue}\bfseries},
    keywords=[2]{plugin, view, page, file, text, config},
    keywordstyle=[2]{\color{blue}\bfseries},
    keywords=[3]{EXT},
    keywordstyle=[3]{\color{blue}\bfseries},
    sensitive=true,
    comment=[l]{\#\ }
}

% PHP5
\lstdefinelanguage[5]{PHP}[]{PHP} {
    morekeywords={
        class,static,private,public,abstract,interface,const,function,require_once,final,new,extends,implements,
        array_combine,array_diff_uassoc,array_udiff,array_udiff_assoc,
        array_udiff_uassoc,array_walk_recursive,array_uintersect_assoc,
        array_usintersect_uassoc,array_uintersect,
        str_split, strpbrk,substr_compare,
        idate,date_sunset,date_sunrise,time_nanosleep,
        return
    }
}

% Fluid
\lstdefinelanguage[Fluid]{HTML}[]{HTML} {
    morekeywords={if, for, each, as, condition, controller, arguments, image, link, action, class}
}

%%%%%%%%%%%%%%%%%%%%%%%%%%%%%%%%%%%%%%%%%%%%%%%%%%%%%%%%%%%%%%%%%%%%%%%
% ZUSÄTZLICHE FUNKTIONALITÄT
%%%%%%%%%%%%%%%%%%%%%%%%%%%%%%%%%%%%%%%%%%%%%%%%%%%%%%%%%%%%%%%%%%%%%%%
% PDF-Einbindung
\usepackage[final]{pdfpages}

% Textformatierung
\usepackage{ragged2e}
\usepackage{lscape}

% Hilfspakete
\usepackage{ifthen}
\usepackage{forloop}
\usepackage{todonotes}
\usepackage{xspace}
\usepackage{datatool}
% !TEX encoding = UTF-8 Unicode
% !TEX root =  Main.tex

% Zeilenabstand 1,5 Zeilen ---------------------------------------------------------------------
\onehalfspacing{}


% Seitenränder ------------------------------------------------------------------------------------
\setlength{\topskip}{\ht\strutbox} % behebt Warnung von geometry
\geometry{paper=a4paper,left=21mm,right=30mm,top=30mm,bottom=20mm}


% Kopf- und Fußzeilen --------------------------------------------------------------------------
\pagestyle{scrheadings}
\renewcommand*{\chapterpagestyle}{scrheadings} 

%% Kopfzeile 
\ihead{} %linke Seite
\chead{} %mittlere Seite
\ohead{\pagemark} %rechte Seite
% \KOMAoptions{headsepline=0.4pt:0.98\paperwidth} % Trennlinie unter Kopfzeile


% Fußzeile
\ifoot{}
\cfoot{} %mittlere Seite
\ofoot{} %rechte Seite
\renewcommand{\footnotesize}{\fontsize{10pt}{12pt}\selectfont} % Fußnoten in 10pt


% sonstige typographische Einstellungen ---------------------------------------------------
\frenchspacing{}
\clubpenalty= 10000
\widowpenalty= 10000 
\displaywidowpenalty= 10000

% Quellcode-Ausgabe formatieren
\lstset{numbers=left, numberstyle=\tiny, numbersep=5pt, breaklines=true}
\lstset{emph={square}, emphstyle=\color{red}, emph={[2]root,base}, emphstyle={[2]\color{blue}}}

% Fußnoten fortlaufend durchnummerieren
\counterwithout{footnote}{chapter}

% Abstand nach Überschriften
\RedeclareSectionCommand[beforeskip=10pt, afterskip=10pt]{chapter}
\RedeclareSectionCommand[beforeskip=10pt, afterskip=5pt]{section}

% Abstand zwischen Kopfzeile und Text
\setlength{\headsep}{15pt} % Standard ist oft 20pt, kleiner machen für weniger Abstand



\begin{document}
% !TEX encoding = UTF-8 Unicode
% !TEX root =  Main.tex

% Trennvorschläge im Text werden mit \" angegeben
% untrennbare Wörter und Ausnahmen von der normalen Trennung können in dieser
% Datei mittels \hyphenation definiert werden

\hyphenation{End-an-wen-der}
\hyphenation{sym-bo-li-siert}
\hyphenation{for-ma-te}
\hyphenation{Ma-nage-ment}
\hyphenation{An-for-de-rungs-be-schrei-bung}
\hyphenation{Cor-po-rate}
\hyphenation{De-sign}
\hyphenation{Pro-gram-mier-schnitt-stel-le}
% !TEX encoding = UTF-8 Unicode
% !TEX root =  Main.tex
% Eigene Befehle und typographische Auszeichnungen für diese

%%%%%%%%%%%%%%%%%%%%%%%%%%%%%%%%%%%%%%%%%%%%%%%%%%%%%%%%%%%%%%%%%%%%%%%
% SCHRIFT UND TYPOGRAPHIE
%%%%%%%%%%%%%%%%%%%%%%%%%%%%%%%%%%%%%%%%%%%%%%%%%%%%%%%%%%%%%%%%%%%%%%%
% Schriftart wechseln
\newcommand{\changefont}[3]{\fontfamily{#1} \fontseries{#2} \fontshape{#3} \selectfont}


%%%%%%%%%%%%%%%%%%%%%%%%%%%%%%%%%%%%%%%%%%%%%%%%%%%%%%%%%%%%%%%%%%%%%%%
% ABKÜRZUNGEN, ZITATE UND REFERENZEN
%%%%%%%%%%%%%%%%%%%%%%%%%%%%%%%%%%%%%%%%%%%%%%%%%%%%%%%%%%%%%%%%%%%%%%%
% Abkürzungen mit korrektem Leerraum
\newcommand{\ua}{\mbox{u.\,a.\ }}
\newcommand{\zB}{\mbox{z.\,B.\ }}
\newcommand{\dahe}{\mbox{d.\,h.\ }}
\newcommand{\Vgl}{Vgl.\ }
\newcommand{\bzw}{bzw.\ }
\newcommand{\evtl}{evtl.\ }

% Hilfsbefehle für Zitate
\newcommand{\maybeAuthor}[1]{%
  \ifblank{\citeauthor{#1}}{o.\,A.}{\citeauthor{#1}}%
}
\newcommand{\maybeYear}[1]{%
  \ifblank{\citeyear{#1}}{o.\,J.}{\citeyear{#1}}%
}

% Zitierformate
\newcommand{\Zitat}[2]{%
  [\maybeAuthor{#1}~\maybeYear{#1}%
  \ifstrempty{#2}{}{, S.~#2}%
  ]%
}
\newcommand{\VglZitat}[2]{%
  [Vgl.\ \maybeAuthor{#1}~\maybeYear{#1}%
  \ifstrempty{#2}{}{, S.~#2}%
  ]%
}
\newcommand{\NarrativZitat}[2]{%
  [\maybeYear{#1}%
  \ifstrempty{#2}{}{, S.~#2}%
  ]%
}
\newcommand{\AutorName}[1]{\textsc{#1}}
\newcommand{\Autor}[1]{\AutorName{\citeauthor{#1}}}

% Abbildungsreferenzen
\newcommand{\Abbildung}[1]{Abbildung~\ref{fig:#1}}
\newcommand{\bs}{$\backslash$}

% verschiedene Befehle um Wörter semantisch auszuzeichnen 
\newcommand{\Fachbegriff}[2][\empty]{\ifthenelse{\equal{#1}{\empty}}{\textit{#2}\xspace}{\textit{#2}\xspace\footnote{#1}\nomenclature{#2}{#1}}}
\newcommand{\FachbegriffSpezialA}[4]{\textit{#4}\footnote{#3}\label{fn:#1}\nomenclature{#4}{#2. Siehe auch Fu{\ss}zeile auf Seite~\pageref{fn:#1}.}}
\newcommand{\FachbegriffSpezialB}[5]{\textit{#5}\footnote{#3}\label{fn:#1}\nomenclature{#4}{#2. Siehe auch Fu{\ss}zeile auf Seite~\pageref{fn:#1}.}}

%%%%%%%%%%%%%%%%%%%%%%%%%%%%%%%%%%%%%%%%%%%%%%%%%%%%%%%%%%%%%%%%%%%%%%%
% WÄHRUNGEN UND BETRÄGE
%%%%%%%%%%%%%%%%%%%%%%%%%%%%%%%%%%%%%%%%%%%%%%%%%%%%%%%%%%%%%%%%%%%%%%%
% Beträge mit verschiedenen Währungssymbolen
\newcommand{\Betrag}[2][general]{#2\,\ifthenelse{\equal{#1}{dollar}}{\$}{}\ifthenelse{\equal{#1}{euro}}{€}{}\ifthenelse{\equal{#1}{yen}}{¥}{}\ifthenelse{\equal{#1}{cent}}{¢}{}\ifthenelse{\equal{#1}{pound}}{£}{}\ifthenelse{\equal{#1}{peso}}{₱}{}\ifthenelse{\equal{#1}{baht}}{฿}{}\ifthenelse{\equal{#1}{franc}}{₣}{}\ifthenelse{\equal{#1}{lira}}{₤}{}\ifthenelse{\equal{#1}{drachma}}{₯}{}\ifthenelse{\equal{#1}{pfennig}}{₰}{}\ifthenelse{\equal{#1}{general}}{¤}{}}

%%%%%%%%%%%%%%%%%%%%%%%%%%%%%%%%%%%%%%%%%%%%%%%%%%%%%%%%%%%%%%%%%%%%%%%
% CODE UND TECHNISCHE AUSZEICHNUNGEN
%%%%%%%%%%%%%%%%%%%%%%%%%%%%%%%%%%%%%%%%%%%%%%%%%%%%%%%%%%%%%%%%%%%%%%%
% Technische Elemente
\newcommand{\Eingabe}[1]{\texttt{#1}}
\newcommand{\Code}[1]{\texttt{#1}}
\newcommand{\Datei}[1]{\texttt{#1}}
\newcommand{\Datentyp}[1]{\textsf{#1}}
\newcommand{\XMLElement}[1]{\textsf{#1}}
\newcommand{\Webservice}[1]{\textsf{#1}}

%%%%%%%%%%%%%%%%%%%%%%%%%%%%%%%%%%%%%%%%%%%%%%%%%%%%%%%%%%%%%%%%%%%%%%%
% TABELLEN UND ABBILDUNGEN
%%%%%%%%%%%%%%%%%%%%%%%%%%%%%%%%%%%%%%%%%%%%%%%%%%%%%%%%%%%%%%%%%%%%%%%
% Beschriftungen anpassen
\addto\captionsngerman{
	\renewcommand{\figurename}{Abb.}
	\renewcommand{\tablename}{Tab.}
}

% Spezielle Spaltentypen
\newcolumntype{w}[1]{>{\raggedleft\hspace{0pt}}p{#1}}
\newcolumntype{y}[1]{>{\RaggedRight\arraybackslash\hsize=#1\hsize}X}



\setcounter{secnumdepth}{3}
\setcounter{tocdepth}{2}


% Deckblatt und Abstract ohne Seitenzahl ---------------------------------------------------------------
% !TEX encoding = UTF-8 Unicode
% !TEX root =  Main.tex

% Kopfbereich mit Logo
\titlehead{%
  % Verschiebt das Logo nach oben
  \vspace*{-1,5cm}%
  % Zentriert das Logo
  \centering%
  % Fügt das HWR-Logo ein (60% der Textbreite)
  \includegraphics[width=0.6\linewidth]{Bilder/Logo_HWR.png}%
  % Abstand nach dem Logo
  \vspace{1,5cm}%  % Weniger Abstand nach dem Logo
}

% Art der Arbeit als Überschrift
\vspace{10cm} 
\subject{\art}

% Erstellt den Titelbereich mit Rahmen
\title{%
  \vspace*{-1,5cm}%  % Verschiebt den gesamten Titelbereich nach oben
  \normalfont\endgraf
  \rule{\textwidth}{.4pt}\\\relax
  \vspace{0.5cm}%  % Kleiner Abstand nach dem Zeilenumbruch
  \begingroup
    \centering
    % Zeilenabstand auf 1.5
    \linespread{1.5}\selectfont
    % Titel in großer Schrift
    \large\titel\\
    % Untertitel in normaler Schrift (falls vorhanden)
    % \normalsize\untertitel
  \endgroup%  % Weniger Abstand vor der unteren Linie
  \endgraf\rule{\textwidth}{.4pt}%
}

% Setzt das Datum und die Hochschule
\date{\vspace{-1.5cm}%  % Reduziert den Abstand zum Titelbereich
\normalsize vorgelegt am \datum\\
an der \\ 
Hochschule für Wirtschaft und Recht Berlin\vspace{5mm}}

% Erstellt die Tabelle mit den Informationen
\publishers{%
  \begin{tabular}{l l}
    % Persönliche Informationen
    \textbf{\normalsize{Von: }} & \normalsize{\autor} \\
    \textbf{\normalsize{Fachrichtung:}} & \normalsize{\fachrichtung} \\
    \textbf{\normalsize{Studienjahrgang:}} & \normalsize{\studienjahrgang} \\
    \textbf{\normalsize{Semester:}} & \normalsize{\semester} \\
    \textbf{\normalsize{Ausbildungsbetrieb:}} & \normalsize{\ausbildungsbetrieb} \\
    % Betreuer-Informationen
    \textbf{\normalsize{Betreuender Prüfer:}} & \normalsize{\erstgutachter} \\
    \textbf{\normalsize{Betreuer Betrieb:}} & \normalsize{\zweitgutachter} \\
    % Unterschriftszeile
    \textbf{\normalsize{Unterschrift Betreuer: }} & \underline{\hspace{3cm}} \\
  \end{tabular}
}

% Keine Seitenzahl auf dem Deckblatt
\pagestyle{empty}
% Erstellt das Deckblatt
\maketitle

% Seitennummerierung -----------------------------------------------------------------------------------
%   Vor dem Hauptteil werden die Seiten in großen römischen Ziffern 
%   nummeriert.
% ------------------------------------------------------------------------------------------------------
\pagenumbering{Roman}
\setcounter{page}{2}
\tableofcontents{}



% Abkürzungsverzeichnis --------------------------------------------------------------------------------
% !TEX encoding = UTF-8 Unicode
% !TEX root =  ../Bachelorarbeit.tex
\chapter*{Abkürzungsverzeichnis}
\markboth{Abkürzungsverzeichnis}{Abkürzungsverzeichnis}
\addcontentsline{toc}{chapter}{Abkürzungsverzeichnis}

\printglossary[type=\acronymtype, title=Abkürzungsverzeichnis]


\newacronym{rup}{RUP}{Rational Unified Process}
\nomenclature{UKW}{Ultrakurzwelle -- Synonym f\"ur UKW-Rundfunk im Bereich 87,5 bis 108 MHz des VHF-Bandes II}

% arabische Seitenzahlen im Hauptteil ------------------------------------------------------------------
\clearpage{}
\pagenumbering{arabic}

% die Inhaltskapitel werden in "Inhalt.tex" inkludiert -------------------------------------------------
% !TEX encoding = UTF-8 Unicode
% !TEX root =  Main.tex

% Hier können die einzelnen Kapitel inkludiert werden. Sie müssen in den 
% entsprechenden .TEX-Dateien vorliegen. Die Dateinamen können natürlich 
% angepasst werden.

% !TEX encoding = UTF-8 Unicode
% !TEX root =  ../Main.tex

\chapter{Einleitung}
\label{cha:Einleitung}


\section{Das Ziel dieser Arbeit}
\label{sec:ZielDerArbeit}
das ist das Ziel 
djufhydkjfhydiofgjydfgomökhhj \gls{rup}

\section{Die Umgebung, in der die Arbeit entstand}
\label{sec:EntstehungsUmgebungArbeit}

Das ist die Umgebung der Arbeit

\section{Der Aufbau dieser Arbeit}
\label{sec:AufbauDieserArbeit}



\begin{description}

	\item[Aktueller Wissensstand:] Der aktuelle Wissensstand beschreibt, auf welchem Wissensniveau sich der Autor im Moment der Aufnahme der Arbeit befand.
	
	\item[Entwicklungsstand TYPO3:] Dieses Kapitel befasst sich mit dem grunds\"atzlichen Entwicklungsstand von TYPO3 Version 4 und 5 und den mit der Extension-Entwicklung zusammenh\"angenden Frameworks Extbase, Fluid und FLOW3. Es werden grundlegende Eigenschaften der Frameworks und deren Leistungsf\"ahigkeit skizziert.
	
	\item[Die Planung des Webradio-Players:] Dieses Kapitel umfasst die Dokumentation der gesamten Planungsphase des Webradio-Players. Hier wird eine Übersicht über die bereits vorhandene Lösung geschaffen und anschließend die zur Planung erforderlichen Dokumente des RUP angefertigt.
	
	\item[Die Entwicklung des Webradio-Players:] Dieses Kapitel enthält die Dokumentation der tatsächlichen Programmierung der Software. Hier werden die Voraussetzungen zur Implementation geklärt und der Verlauf der Entwicklung anhand von Beispielen schrittweise abgearbeitet.
	
	\item[Fazit und kritische Bewertung:] Im Fazit werden die gemachten Erfahrungen und die Ergebnisse der Planung und Entwicklung abschließend zusammengefasst und kritisch bewertet. Zusätzlich wird ein kleiner Ausblick auf Erweiterungsmöglichkeiten und mögliche Optimierungsschritte unternommen.

\end{description}


\section{Wenige Informationen, wenige Quellen \dots}
\label{sec:Quellenlage}
das ist was anders

%\include{Inhalt/Wissensstand}

%\include{Inhalt/Entwicklungsstand}

%\include{Inhalt/Methoden}

%\include{Inhalt/Planung/Planung}
	%\include{Inhalt/Planung/Vision}
	%\include{Inhalt/Planung/Anforderungsbeschreibung}
	%\include{Inhalt/Planung/Architektur}
	
	
%\include{Inhalt/Entwicklung/Entwicklung}
	%\include{Inhalt/Entwicklung/ExtensionBuilder}
	%\include{Inhalt/Entwicklung/BeginnProgrammierung}
	%\include{Inhalt/Entwicklung/Implementierung}

% !TEX encoding = UTF-8 Unicode
% !TEX root =  ../Bachelorarbeit.tex
\chapter{Fazit und kritische Bewertung}
\label{cha:Fazit}


\section{Das Ergebnis}
\label{sec:Ergebnis}

Der Relaunch des gesamten Webauftritts von HIT RADIO FFH erfolgte am 24.\,Juli 2011, mit dem auch der neue Webradio-Player zum ersten Produktiveinsatz gelangte. In Abbildung~\ref{fig:FfhWebradioEndversion} ist die Webradio-Player-Extension im Live-Einsatz auf \href{http://webradio.ffh.de/}{webradio.ffh.de} zu sehen. 

%% Screenshot der Endversion
\begin{figure}[htbp]
	\begin{center}
		\includegraphics[width=\textwidth]{Fazit/FfhWebradioEndversion.png}
		\caption[FFH Webradio-Player im Live-Einsatz]{Der FFH Webradio-Player im Live-Einsatz}
		\label{fig:FfhWebradioEndversion}
	\end{center}
\end{figure}

Die grunds\"atzliche Entscheidung, bei der Extension-Entwicklung auf die Frame\-works Extbase und Fluid zu setzen, hat sich als absolut richtig erwiesen. Trotz der vielen Kinderkrankheiten und der unvollst\"andigen Implementation vieler Funktionen ist die anwendungsdom\"anen-getriebene Herangehensweise von Extbase bzw. FLOW3 ein deutlicher Schritt in eine einfachere Richtung der Entwicklung. 

Der Webradio-Player l\"auft seit dem Relaunch nahezu fehlerfrei und ben\"otigte danach nur kleine Anpassungen wegen Schnittstellen\"anderungen des Stream-Providers Nacamar.


\newpage

\section{Die Bewertung der Frameworks Extbase und Fluid}
\label{sec:BewertungFrameworks}


Das bei der Radio/Tele FFH GmbH im Zuge des Relaunches zum Einsatz kommende Content Management System TYPO3 bildete die Basis der Extension-Entwicklung w\"ahrend dieser Arbeit. Aufgrund der angepriesenen Zukunftssicherheit wurde zudem entschieden, alle Extensions mit denen am Markt neu eingef\"uhrten Frameworks Extbase und Fluid zu entwickeln. Sowohl Extbase als MVC und Domain-Driven Design basiertes PHP-Framework, als auch Fluid als XML basierte Templating-Engine vereinen das gemeinsame Prinzip »Convention over Configuration«. Dieses Prinzip stellt viele Entwickler vor die gro{\ss}e Herausforderung, ihre bisherigen Gewohnheiten beim Entwickeln von Software zu \"uberdenken.

Die vielen zun\"achst \"ubertrieben erscheinenden Konventionen beim Programmieren mit Extbase und Fluid haben schlussendlich trotz vieler Zweifel erm\"oglicht, dass drei Entwickler parallel an verschiedenen Extensions arbeiten konnten, auch wenn sie nicht am urspr\"unglichen Planungsprozess beteiligt waren. Einzig durch die strengen Konventionen in der Programmierung gelingt es, sich ohne gro{\ss}e Probleme in andere Extensions einzulesen. So war es beispielsweise m\"oglich, bestimmte Quellcode-Fragmente schnell in andere Extensions zu portieren, ohne viel am bereits vorhandenen Code \"andern zu m\"ussen.

Die Templating-Engine Fluid \"ubernimmt diese hervorragenden Eigenschaften bei der Umsetzung in der View-Ebene der Extensions. Als XML-basierte Templating-Engine ist sie flexibel einsetzbar, erzeugt validen XHTML-Code und ist durch ihre Viewhelper theoretisch unendlich erweiterbar. Die bereits implementierten Viewhelper zeugen von der Macht, die von dieser Engine ausgeht. Aus diesem Grund ist es schade, dass bisher nur Basisfunktionen in Viewhelpern vorliegen. Geht es darum, exotischere Logiken in Viewhelpern zu verwenden, so ist man leider darauf angewiesen, diese selbst zu programmieren. Hierdurch geht ein nicht unerheblicher Teil an Entwicklungszeit verloren, weil die Rahmenbedingungen f\"ur die Umsetzung TYPO3 interner Funktionen leider meistens schlecht sind. Es fehlt an Dokumentationen der TYPO3 Funktionalit\"aten, aber auch an M\"oglichkeiten, diese in Viewhelpern zu realisieren. Aus diesem Grund ist man auch hier stark auf die Community im Internet angewiesen. Den Hauptanteil an Informationen bezieht man von anderen Entwicklern, die sich die L\"osungen meist umst\"andlich durch \Fachbegriff[Heuristische Methode, bei der durch Versuch und Irrtum eine L\"osung gefunden wird.]{trial \& error} erarbeitet haben.


Theoretisch l\"asst sich Fluid mit jedem verf\"ugbaren PHP-Framework verwenden. So kann eine simple Webseite auch komplett autark ohne Content Management System aufgebaut werden.


Zusammenfassend bleibt zu sagen, dass sowohl Extbase als auch Fluid ihren Weg in die Welt von TYPO3 gefunden haben. Die beiden alternativen Frameworks geben einen Vorgeschmack auf die Entwicklung mit TYPO3 v5 und FLOW3 und helfen Entwicklern schon heute sich auf die Portierung ihrer Extensions vorzubereiten. F\"ur den Autor als Neueinsteiger in der Extension-Entwicklung boten die beiden Werkzeuge viele neue und intuitive Wege, an die Entwicklung von Software heranzugehen. Nicht zuletzt ist hervorzuheben, dass die neuartigen Ans\"atze der Entwicklung, wie etwa das Domain-Driven Design und Convention over Configuration, viel Erleichterung in den Alltag eines Software-Entwicklers bringen k\"onnen, sofern dieser sich darauf einl\"asst.



\section{Ein Ausblick}
\label{sec:EinAusblick}

Es hat sich gezeigt, dass die Entwicklung der Frameworks Extbase und Fluid st\"andig voranschreitet, weshalb es in unregelm\"a{\ss}igen Abst\"anden sinnvoll ist, die Webradio-Player-Extension zu aktualisieren und Neuerungen direkt zu implementieren. Dadurch ist eine sp\"atere Umsetzung in FLOW3 f\"ur TYPO3 Version 5 einfacher zu realisieren. 

Bereits im Verlaufe der Erstellung dieser Arbeit ergaben sich \"Anderungen in Extbase, die dazu f\"uhrten, dass die Methode zum Instanziieren von Repository-Objekten ge\"andert wurde. Fr\"uher geschah dies \"uber eine \FachbegriffSpezial{Factory}{Entwurfsmuster, wirkt als Hilfsmittel zur Erzeugung von Produkten und Objekten}{»Eine Factory ist ein Hilfsmittel zur Erzeugung von Produkten und Objekten, unabhängig von konkreten Klassen. Sie wird verwendet, wenn die Instanziierung nur über den new()-Operator nicht möglich oder nicht sinnvoll ist, zum Beispiel wenn das Objekt schwierig zu konstruieren ist oder vorher konfiguriert werden muss. [\dots] Factory-Methoden sind statische Methoden zur Erzeugung von Objekten des eigenen Klassentyps. Ein Singleton kann zum Beispiel als Spezialfall einer Klasse mit einer Factory-Methode angesehen werden.« \ZitatSpezial{Vgl.}{Abschnitt „Factory“}{horn:Patterns}}{Factory}-Methode der TYPO3 Bibliothek \Code{t3lib\_div}. Mittlerweile wird ein Repository \"uber eine \FachbegriffSpezial{DependencyInjection}{Entwurfsmuster, soll helfen Abh\"angigkeiten eines Objektes einfacher aufzul\"osen}{»Dependency Injection ist ein Entwurfsmuster (Design Pattern), dass helfen soll, Abhängigkeiten eines Objektes in der objekt-orientierten Programmierung einfacher aufzulösen. Es findet eine Umkehr der Steuerung (Inversion of Control) statt, um das Objekt von unnötigen Verbindungen zu seiner Umwelt zu befreien, die es nur für die Auflösung von Abhängigkeiten, nicht aber für seine eigentliche Aufgabe benötigt. Die Verantwortung für das Auflösen der Abhängigkeiten wird aus dem Objekt in das umliegende Framework, in unserem Falle Extbase, verlagert.« \Zitat[]{oertel:DependencyInjection}}{Extbase Dependency Injection} erzeugt.

Durch die grundverschiedenen Ans\"atze beider Varianten ist schnell erkennbar, dass mit der Entwicklung Schritt gehalten werden muss, damit die eigenen Extensions zukunftssicher bleiben. Aus diesem Grund werden alle bei der Radio/Tele FFH GmbH entwickelten Extensions -- so auch der Webradio-Player -- st\"andig auf dem neusten Entwicklungsstand gehalten. 

Einer Portierung auf FLOW3 st\"unde somit nichts im Wege.







% Literaturverzeichnis ---------------------------------------------------------------------------------
%   Das Literaturverzeichnis wird aus der BibTeX-Datenbank "Bibliographie.bib"
%   erstellt.
% ------------------------------------------------------------------------------------------------------
\bibliography{Literatur/Bibliographie} % Aufruf: bibtex Masterarbeit
\bibliographystyle{natdin} % DIN-Stil des Literaturverzeichnisses


% Restliche Verzeichnisse -> wenn diese leer sind, sollte das hier auskommentiert werden ------------------------------------------------------------------------------
\listoffigures{} % Abbildungsverzeichnis
\listoftables{} % Tabellenverzeichnis


% Ehrenwörtliche Erklärung ----------------------------------------------------------------------------
% !TEX encoding = UTF-8 Unicode
% !TEX root =  Main.tex

%muss nicht mehr geändert werden, ist schon nach HWR Vorlage

\addchap{Ehrenwörtliche Erklärung}
Ich erkläre ehrenwörtlich:

dass ich die vorliegende \art\xspace in allen Teilen selbstständig angefertigt und keine anderen als die in der Arbeit angegebenen Quellen und Hilfsmittel benutzt habe, und dass die Arbeit in gleicher oder ähnlicher Form in noch keiner anderen Prüfung vorgelegen hat. Sämtliche wörtlichen oder sinngemäßen Übernahmen und Zitate, sowie alle Abschnitte, die mithilfe von Kl- basierten Tools entworfen, verfasst und/oder bearbeitet wurden, sind kenntlich gemacht und nachgewiesen. Im Anhang meiner Arbeit habe ich sämtliche KI-basierte Hilfsmittel angegeben. Diese sind mit Produktnamen und formulierten Eingaben (Prompts) in einem Kl-Verzeichnis ausgewiesen.

Ich bin mir bewusst, dass die Verwendung von Texten oder anderen Inhalten und Produk- ten, die durch Kl-basierte Tools generiert wurden, keine Garantie für deren Qualität darstellt. Ich verantworte die Ubernahme jeglicher von mir verwendeter maschinell generierter Passagen vollumfänglich selbst und trage die Verantwortung für eventuell durch die Kl generierte fehlerhafte oder verzerrte Inhalte, fehlerhafte Referenzen, Verstöße gegen das Datenschutz- und Urheberrecht oder Plagiate.\\

\ort, den \today

% hier kann eine Unterschrift eingefügt werden, wenn die Arbeit ausgedruckt wird
% \begin{flushright}
%     \includegraphics[height=2cm]{Bilder/Unterschrift.png}
% \end{flushright}


\autor\\
 


% Anhang -----------------------------------------------------------------------------------------------
%   Die Inhalte des Anhangs werden analog zu den Kapiteln inkludiert.
%   Dies geschieht in der Datei "Anhang.tex".
% ------------------------------------------------------------------------------------------------------
\appendix
    % \clearpage{}
    % \pagenumbering{roman}
    \chapter{Anhang}
    \label{sec:Anhang}
    % Rand der Aufzählungen in Tabellen anpassen
    \setdefaultleftmargin{1em}{}{}{}{}{}
    % !TEX encoding = UTF-8 Unicode
% !TEX root =  Main.tex

% Beispiel: 
% \section{Beiliegende CD}
% \label{sec:BeiliegendeCd}



\end{document}
